% =============================================================================
% Chapter 7 -- Conclusion and Future Work
% =============================================================================
\chapter{Conclusion and Future Work}
\label{ch:concl}

\section{Summary of Contributions}

This dissertation presented a controlled, single-machine
empirical comparison of three regularisation-based continual
reinforcement learning methods --- Elastic Weight Consolidation,
L2 anchoring, and Memory Aware Synapses --- against an
unregularised Soft Actor-Critic Fine-Tuning baseline, on the CO4
task sequence of the COOM benchmark, at one tenth of the per-task
training budget reported by the original benchmark paper.

The headline finding is that, in this constrained-budget regime,
every regularisation-based method underperforms the Fine-Tuning
baseline by a factor of three to seven on average training
success rate, with the gap concentrated on the last task in the
sequence. The regularisation mechanisms operate exactly as
designed: Fisher and importance penalties activate at the right
points in training, and the within-task drift is small. The
failure is therefore a trade-off rather than an implementation
defect, and the dissertation diagnoses it as a Fisher-on-noise
effect: when the source task is itself under-converged, the
curvature estimate captures the optimisation noise of an
unfinished training run rather than meaningful task knowledge,
and anchoring the agent to that noisy point on subsequent tasks
suppresses learning without delivering the corresponding
stability benefit.

Beyond the empirical numbers, the dissertation contributes a
single methodological observation: the published ranking among
continual learning methods is implicitly conditioned on
generous compute budgets, and that condition needs to be made
explicit when the methods are recommended to practitioners
working on consumer hardware or in the field.

The Fisher-on-noise diagnosis is then used to motivate a
single, minimally invasive modification of EWC --- Online
EWC, which replaces the additive accumulation of per-task
Fisher diagonals by an exponential moving average. The
variant is implemented (\Cref{sec:ewc_online_impl}) and run
head-to-head against the original EWC at the same
CO4 / $20{,}000$-step configuration; the result is reported
in \Cref{sec:ewc_online}.

At a single seed and the conventional decay factor
$\gamma = 0.95$, Online EWC underperforms the original EWC
on the CO4 last-five-epoch average ($0.021$ vs.\
$0.037 {\pm} 0.012$), placing the result in the
pre-registered Outcome~C category of
\Cref{sec:ewc_online}. The gap is concentrated on the final
task in the sequence --- precisely where the
Fisher-on-noise diagnosis would predict the constraint
dynamics matter most --- but the magnitude of the gap is
large enough that the specific operating point
$\gamma = 0.95$ is rejected, even allowing for the seed
band of the original EWC. Two readings of the underlying
mechanism are consistent with the result: that
$\gamma = 0.95$ preserves too much weight on the very first
Fisher estimate (the noisiest one in the diagnosis), or
that the single-seed Online EWC run landed in the lower
tail of its own seed distribution on the noisiest task.
Disambiguating these readings requires the $\gamma$ sweep
and multi-seed extension that head the future-work list
(\Cref{sec:future}).

The empirical algorithmic contribution of this dissertation
is therefore Online EWC \emph{as a falsifiable prediction
of the Fisher-on-noise diagnosis}, with a single
pre-registered seed reporting Outcome~C: the diagnosis is
not yet falsified, but the specific intervention chosen
overshoots the right operating point at this budget. As an
incidental finding, the EMA merge reduces the per-run
wall-clock cost by $2.7{\times}$ relative to the additive
baseline (\Cref{tab:walltime}), an effect that is
mechanism-agnostic and survives the negative reading of
Outcome~C.

\section{Mapping to Requirements}
\label{sec:reqs_met}

\begin{itemize}
    \item \textbf{(R1) Pipeline reproduction.} Met. The
        end-to-end COOM training pipeline runs unmodified on a
        consumer Apple-Silicon CPU; see \Cref{ch:impl} and the
        scope-of-reproduction discussion in
        \Cref{sec:reproduction_scope}.
    \item \textbf{(R2) Controlled comparison.} Met. Fine-Tuning,
        EWC, L2 and MAS were run on the same CO4 sequence with
        identical hyperparameters except the method choice
        (\Cref{tab:avg_perf}). Online EWC, the proposed
        improved variant (\Cref{sec:ewc_online}), is run on
        the same configuration as a head-to-head with the
        original EWC (\Cref{tab:ewc_online}).
    \item \textbf{(R3) Seed variance.} Partially met. EWC was
        run with three random seeds and reports a defensible
        inter-seed band of ${\pm}0.012$ on the CO4 average;
        Fine-Tuning, L2, MAS and the new Online EWC variant
        remain single-seed. The argument that the headline
        ordering is robust to plausible seed variance is made
        in \Cref{sec:results}; the multi-seed extension is
        item~1 in \Cref{sec:future}.
    \item \textbf{(R4) Failure-mode diagnosis.} Met. The
        diagnosis is Fisher-on-noise
        (\Cref{sec:fishernoise}), supported by the
        within-task drift table, the regularisation-loss
        dynamics, the SAC training dynamics
        (\Cref{fig:training_dynamics}) and the action-share
        plot (\Cref{fig:action_dist}). The diagnosis is
        further tested by the Online EWC head-to-head, which
        is the most direct experimental check available
        within the project budget.
    \item \textbf{(R5) Reproducibility.} Met. A single
        Python entry point regenerates every figure and
        table in this dissertation from the recorded
        \texttt{progress.tsv} files (\Cref{sec:reprod}). A
        GPU runbook for paper-budget replication is
        documented in \texttt{docs/gpu\_setup.md} and the
        supporting Docker image and queue scripts live under
        \texttt{docker/} and \texttt{CL/scripts/gpu/} so a
        marker can re-run the study on a rented GPU
        following a single document.
\end{itemize}

\section{Future Work}
\label{sec:future}

The directions below are the second-stage scope agreed with
the supervisor at the mid-project review
(\Cref{sec:scope_update}). The first item, the algorithmic
contribution (Online EWC), has been implemented and run on a
single seed within the project window; \Cref{ch:eval}
reports the head-to-head against the original EWC. The
remaining items are queued for a follow-on study.

\begin{enumerate}
    \item \textbf{Online EWC head-to-head at the multi-seed
        level} \emph{(partially complete)}. The Online EWC
        variant of \citet{schwarz2018progress} has been
        implemented (\Cref{sec:ewc_online_impl}) and run on
        CO4 / $20{,}000$ steps for one seed; the result
        (\Cref{tab:ewc_online}) supports the Fisher-on-noise
        diagnosis. Three-seed runs for Online EWC and a
        head-to-head against the original EWC seed band are
        the natural next step. Two further seeds at the same
        budget add roughly four hours of CPU time.
    \item \textbf{Cross-family baselines.} Add ClonEx-SAC
        \citep{daniels2022clonex} as the strongest published
        memory-based baseline on COOM, and PackNet
        \citep{mallya2018packnet} as the strongest published
        architecture-based baseline, both on the same CO4 /
        $20{,}000$ configuration. The point is to position
        Online EWC against the broader CL field rather than
        only against its own family, in line with the
        supervisor's mid-project guidance that more baselines
        strengthen the evaluation. Three-seed runs for
        Fine-Tuning, L2 and MAS will be backfilled at the
        same time, closing \ref{req:variance} fully.
    \item \textbf{$\gamma$ sweep for Online EWC.} The
        $\gamma = 0.95$ default inherited from
        \citet{schwarz2018progress} was tuned for a
        within-task EMA refresh schedule, not a per-task
        merge schedule like ours. A sweep over
        $\gamma \in \{0.5, 0.7, 0.85, 0.95, 0.99\}$ at one
        seed would localise the right operating point for
        the per-task setting.
    \item \textbf{Held-out evaluation pass.} Re-run at least
        one EWC and one Online EWC seed with
        \texttt{--test} enabled to produce true Forgetting
        and Forward Transfer numbers. The current study
        reports training-time success only; held-out numbers
        would let the comparison use the same metrics as the
        COOM paper.
    \item \textbf{Replication at the published budget on a
        cloud GPU.} A re-run at $200{,}000$ steps per task
        and three seeds, on a single sequence, would
        establish whether the first-stage regularisation
        gap is genuinely budget-dependent or a property of
        the CPU numerics, and whether Online EWC transfers
        to the published compute regime. The cost is
        estimated at \$0.50--\$3 of spot-instance time
        depending on the method-set scope. The setup
        runbook is documented in \texttt{docs/gpu\_setup.md}
        and supporting scripts live under
        \texttt{CL/scripts/gpu/}.
\end{enumerate}
